The CR5 stack required me to align HDF5 records, preserve state and action semantics through LeRobot conversion, and serve policy chunks over WebSocket, with schema and motion checks around their use. For the HERO RoboMaster team, I helped modularize image acquisition, vision processing, and communication synchronization as ROS2 nodes. These experiences taught me to make interfaces testable, yet I lack systematic training in the protocols, software lifecycle, and distributed-data design needed for reliable full-stack IT. At The Hong Kong Polytechnic University (PolyU), the MSc in Information Technology would close that gap.

If the current curriculum is offered in my intake, COMP5241 Software Engineering and Development would strengthen how I specify, test, and maintain cross-component contracts, while COMP5311 Internet Infrastructure and Protocols would deepen the networking reasoning behind remote policy serving. Subject to elective availability, COMP5434 Big Data Computing would help me move synchronized robot records beyond a single-machine pipeline. If I chose the optional project route and the six-credit COMP5933 Project were approved, I could build a distributed prototype whose integration tests trace schemas, timestamps, service responses, and execution behavior end to end. My experience locating faults at data and execution boundaries could ground those tests in concrete failure cases, preparing me to develop dependable infrastructure for embodied-agent R\&D.
